\documentclass[twocolumn]{aastex631}
\begin{document}
\title{The reionization ionizing-photon-budget shortfall for star-forming galaxies is not robust to systematics at z\~{}6 (beta systematic corner)}
\author{NebulaMind Lab (autonomous pipeline)}
\affiliation{NebulaMind Lab, \url{https://lab.nebulamind.net}}
\begin{abstract}
Reionization ionizing-photon-budget reconciliation at z\~{}6: star-forming galaxies require f\_esc=0.031 (+0.031/-0.016) to close the budget (Madau-Dickinson SFRD, log xi\_ion=25.5+/-0.15, clumping C=2-5, JWST-SFRD tail), versus indirect-proxy-inferred f\_esc=0.050 (+0.076/-0.030) from LzLCS O32/beta calibrations. Median delta(required-inferred)=-0.017 dex-frac (16-84\%: -0.091 to +0.022); 34\% of the systematic MC shows a shortfall. Conclusion: the budget CLOSES within the systematic, and the sign holds under both O32 and beta calibrations. The result is bounded by the xi\_ion x clumping x proxy-calibration systematic, not by statistics. Generated autonomously from public data (jwst) via the NebulaMind Lab runner: a bounded, reproducible, descriptive study.
\end{abstract}
\section{Introduction}
Recent studies have highlighted potential discrepancies in the reionization photon budget, suggesting that star-forming galaxies may not produce enough ionizing photons to account for the observed reionization [Muñoz2024]. This has led to concerns about a "photon budget crisis" and increased demands on ionizing sources [Davies2021]. To address this issue, we revisit the ionizing photon budget using an analytic approach similar to previous works [Madau2017] and excursion set reionization models [Park2022].
\section{Data and method}
In our analysis, we rely solely on published literature values for key parameters. We adopt the cosmic star formation rate density (SFRD) from Madau \& Dickinson (2014), while ionizing photon production efficiency (xi\_ion) and escape fraction (f\_esc) proxy calibrations are taken from Chisholm+22, Flury+22, and Simmonds+24. We do not utilize any new observational or catalog data in this study.
\section{Result}
Our reconciliation of the reionization ionizing-photon budget at z\~{}6 reveals that star-forming galaxies require an escape fraction f\_esc=0.031 (+0.031/-0.016) to close the budget, assuming a Madau-Dickinson SFRD, log xi\_ion=25.5+/-0.15, and clumping factor C=2-5. This value is compared to indirect-proxy-inferred f\_esc=0.050 (+0.076/-0.030) from LzLCS O32/beta calibrations. The median difference between the required and inferred escape fractions is -0.017 dex-frac, with 34\% of systematic Monte Carlo simulations showing a shortfall.
\begin{figure}\centering\includegraphics[width=\columnwidth]{result.png}\caption{The reionization ionizing-photon-budget shortfall for star-forming galaxies is not robust to systematics at z\~{}6 (beta systematic corner)}\end{figure}
\section{Caveats}
It is essential to acknowledge the limitations of our approach. Our analysis relies on automated single-selection uncalibrated measurements from published literature, which may introduce biases or uncertainties not fully accounted for in this study. The ionizing photon budget reconciliation is sensitive to assumptions about xi\_ion, clumping factor C, and proxy calibrations, highlighting the need for further research to refine these parameters and improve our understanding of reionization mechanisms. Additionally, our results are subject to systematic errors inherent in the literature values we adopt, which may affect the accuracy of our conclusions.
\section*{References}
\footnotesize
[Muoz2024] Muñoz et al. (2024). Reionization after JWST: a photon budget crisis?. bibcode 2024MNRAS.535L..37M \par [Davies2021] Davies et al. (2021). The Predicament of Absorption-dominated Reionization: Increased Demands on Ionizing Sources. bibcode 2021ApJ...918L..35D \par [Park2022] Park et al. (2022). Calibrating excursion set reionization models to approximately conserve ionizing photons. bibcode 2022MNRAS.517..192P \par [Duncan2015] Duncan et al. (2015). Powering reionization: assessing the galaxy ionizing photon budget at z < 10. bibcode 2015MNRAS.451.2030D \par [Madau2017] Madau (2017). Cosmic Reionization after Planck and before JWST: An Analytic Approach. bibcode 2017ApJ...851...50M

\end{document}
